% ===================================================================
%  Shared Beamer preamble for "Practical Image Processing"
%  English slides. Compile with pdfLaTeX.
% ===================================================================

\usetheme{metropolis}
\usepackage[T1]{fontenc}
\usepackage[utf8]{inputenc}
\usepackage{lmodern}
\usepackage{listings}
\usepackage{tcolorbox}
\usepackage{fontawesome5}
\usepackage{booktabs}
\usepackage{amsmath,amssymb}
\usepackage{tikz}
\usepackage{pgfplots}
\pgfplotsset{compat=1.17}
\usetikzlibrary{arrows.meta,positioning,shapes.geometric,fit,backgrounds,matrix}
\tcbuselibrary{skins,breakable}

% ---- Palette --------------------------------------------------------
\definecolor{coursePrimary}{HTML}{1B4965}   % deep blue
\definecolor{courseAccent}{HTML}{E76F51}    % coral
\definecolor{courseGood}{HTML}{2A9D8F}      % teal
\definecolor{courseBad}{HTML}{C0392B}       % red
\definecolor{courseInfo}{HTML}{2980B9}      % blue
\definecolor{codeBg}{HTML}{F5F6FA}
\definecolor{codeKw}{HTML}{8E44AD}
\definecolor{codeStr}{HTML}{2A9D8F}
\definecolor{codeCom}{HTML}{7F8C8D}

\setbeamercolor{frametitle}{bg=coursePrimary}
\setbeamercolor{progress bar}{fg=courseAccent}
\setbeamercolor{alerted text}{fg=courseAccent}
\metroset{block=fill,sectionpage=progressbar,numbering=fraction}

% ---- Code listings (Python is built into listings) -----------------
\lstdefinestyle{course}{
  backgroundcolor=\color{codeBg},
  basicstyle=\ttfamily\scriptsize,
  keywordstyle=\color{codeKw}\bfseries,
  stringstyle=\color{codeStr},
  commentstyle=\color{codeCom}\itshape,
  numberstyle=\tiny\color{codeCom},
  numbers=left, numbersep=6pt,
  showstringspaces=false, breaklines=true, breakatwhitespace=true,
  tabsize=4, columns=fullflexible,
  frame=single, framesep=4pt, rulecolor=\color{gray!30},
  keepspaces=true, upquote=true, language=Python,
}
\lstset{style=course}
\lstdefinelanguage{bash}{
  keywords={cd,pip,python,venv,source,import}, sensitive=true,
  morecomment=[l]{\#}, morestring=[b]', morestring=[b]",
}

% ---- Reusable callout boxes ----------------------------------------
\newtcolorbox{goodbox}[1][Good practice]{colback=courseGood!8,colframe=courseGood,
  fonttitle=\bfseries,title={\faCheck~#1},boxrule=0.6pt,left=4pt,right=4pt,top=3pt,bottom=3pt}
\newtcolorbox{badbox}[1][Pitfall]{colback=courseBad!8,colframe=courseBad,
  fonttitle=\bfseries,title={\faTimes~#1},boxrule=0.6pt,left=4pt,right=4pt,top=3pt,bottom=3pt}
\newtcolorbox{infobox}[1][Note]{colback=courseInfo!8,colframe=courseInfo,
  fonttitle=\bfseries,title={\faInfoCircle~#1},boxrule=0.6pt,left=4pt,right=4pt,top=3pt,bottom=3pt}
\newtcolorbox{mathbox}[1][Math]{colback=coursePrimary!6,colframe=coursePrimary,
  fonttitle=\bfseries,title={\faCalculator~#1},boxrule=0.6pt,left=4pt,right=4pt,top=3pt,bottom=3pt}

\newcommand{\keyidea}[1]{\begin{alertblock}{Key idea}#1\end{alertblock}}
\newcommand{\term}[1]{\textbf{\color{coursePrimary}#1}}
\newcommand{\code}[1]{\lstinline[basicstyle=\ttfamily\small]|#1|}

% ---- A helper to draw a small pixel grid, cells shaded by value ----
% \pixelgrid{ {10,20,120,200,240}, {30,60,130,180,230}, ... }
\newcommand{\pixelgrid}[1]{%
  \begin{tikzpicture}[font=\scriptsize,scale=0.6]
    \foreach \row [count=\ry] in {#1} {
      \foreach \val [count=\cx] in \row {
        \draw[gray,fill=coursePrimary!8] (\cx,-\ry) rectangle ++(1,-1);
        \node at (\cx+0.5,-\ry-0.5) {\val};
      }
    }
  \end{tikzpicture}}
